\begin{abstract}
Text classifiers receive Unicode strings, while readers primarily perceive rendered glyphs. This mismatch enables adversarial changes that preserve the apparent text but alter its code points and tokenization. We compare word and character TF--IDF classifiers, an \texttt{o200k\_base} BPE-feature classifier, and fine-tuned DistilBERT on BANKING77 intent detection and coarse TREC question classification. The evaluation covers homoglyph substitutions, invisible characters, compatibility characters, canonical composition, and mixed attacks at four strengths. We additionally evaluate Unicode normalization, removal of default-ignorable characters, and an English-specific confusable-repair heuristic. Beyond accuracy and Macro-F1, we measure attack success, defense recovery, token-sequence disruption, and the association between tokenization changes and failures. The resulting study separates controlled representation effects from the behavior of a pretrained Transformer and identifies both the value and the limits of low-cost input sanitization.
\end{abstract}

